\usepackage{amsmath,amssymb,bm,mathtools,esint}
\usepackage{graphicx,booktabs,array,tabularx,multicol}
\usepackage{siunitx}
\setCJKmainfont[AutoFakeBold=2]{SimHei}
\setCJKsansfont[AutoFakeBold=2]{SimHei}
\setCJKmonofont[AutoFakeBold=2]{SimHei}
\setsansfont{TeX Gyre Heros}
\usefonttheme{professionalfonts}
\geometry{paperwidth=254mm,paperheight=143mm}
\definecolor{courseblue}{HTML}{2B4D89}
\definecolor{courselight}{HTML}{DFEBF5}
\definecolor{coursepale}{HTML}{7D9ED1}
\definecolor{coursered}{HTML}{B22222}
\setbeamercolor{normal text}{fg=black,bg=white}
\setbeamercolor{structure}{fg=courseblue}
\setbeamercolor{frametitle}{fg=black,bg=white}
\setbeamercolor{block title}{fg=white,bg=courseblue}
\setbeamercolor{block body}{fg=black,bg=courselight}
\setbeamercolor{alerted text}{fg=coursered}
\setbeamercolor{footline}{fg=white,bg=coursepale}
\newcommand{\coursepart}{课程安排}
\setbeamerfont{normal text}{size=\fontsize{17.5}{22}\selectfont}
\setbeamerfont{frametitle}{size=\fontsize{24}{29}\selectfont,series=\bfseries}
\DeclareMathSizes{24}{24.88}{17.28}{12}
\setbeamerfont{block title}{size=\fontsize{17.5}{21}\selectfont,series=\bfseries}
\setbeamerfont{block body}{size=\fontsize{17.5}{22}\selectfont}
\setbeamertemplate{navigation symbols}{}
\setbeamertemplate{blocks}[rounded][shadow=false]
\setbeamertemplate{itemize item}{\color{courseblue}\raisebox{.15ex}{\rule{.36ex}{.36ex}}}
\setbeamertemplate{itemize subitem}{\color{courseblue}\rule{.30ex}{.30ex}}
\setbeamertemplate{frametitle}{%
  \vspace{0.12cm}\centering\insertframetitle\par\vspace{0.16cm}%
  \makebox[\textwidth][c]{\color{courseblue}\rule{\dimexpr\paperwidth-0.7cm\relax}{2.5pt}}\par\vspace{-0.12cm}}
\setbeamertemplate{footline}{%
  \leavevmode\hbox{%
   \begin{beamercolorbox}[wd=.80\paperwidth,ht=15pt,dp=4pt,leftskip=12pt]{block title}%
    \fontsize{12}{14}\selectfont\coursepart
   \end{beamercolorbox}%
   \begin{beamercolorbox}[wd=.20\paperwidth,ht=15pt,dp=4pt,rightskip=12pt plus1fil]{footline}%
    \hfill\fontsize{12}{14}\selectfont\insertframenumber\,/\,\inserttotalframenumber
   \end{beamercolorbox}}}
\setbeamersize{text margin left=0.65cm,text margin right=0.65cm}
\setlength{\parskip}{5pt}
\setlength{\abovedisplayskip}{8pt}
\setlength{\belowdisplayskip}{8pt}
\AtBeginDocument{\fontsize{17.5}{22}\selectfont}
\AtBeginEnvironment{frame}{\fontsize{17.5}{22}\selectfont}
\setlength{\emergencystretch}{1.5em}
\newcommand{\dd}{\mathop{}\!\mathrm d}
\newcommand{\epso}{\varepsilon_0}
\newcommand{\muo}{\mu_0}
\newcommand{\vct}[1]{\bm{#1}}
\renewcommand{\unit}[1]{\,\mathrm{#1}}
\newcommand{\e}{\mathrm e}
\newcommand{\ii}{\mathrm i}
\newcommand{\src}[1]{\note{来源：#1}}
\newcommand{\teach}[1]{\note{讲授：#1}}
\graphicspath{{figures/}}
\hypersetup{pdftitle={基础物理学III：电磁学小班辅导},pdfauthor={小班辅导讲义},colorlinks=true,linkcolor=courseblue,urlcolor=courseblue}
